%! AP Statistics — Unit 9, Lesson 9.1 — Group Activity
\documentclass[10pt]{article}
\usepackage{apstats-article}
\usepackage{apstats-boxes}

\begin{document}

\pageheader{Unit 9, Lesson 9.1}{Group Activity}
\namepartnerperiod

\vspace{0.1in}
\textbf{Is the Slope Real? Exploring Sampling Variability in Regression}

\vspace{0.05in}
Below are three regression scenarios. In each, a researcher collected data and fit a
least-squares regression line. Your job: decide whether the observed slope seems
convincingly different from 0, or whether it could plausibly be due to random variation.

\vspace{0.1in}
\begin{center}
\renewcommand{\arraystretch}{1.4}
\begin{tabular}{clccc}
    \textbf{\#} & \textbf{Context} & \textbf{Explanatory ($x$)} & \textbf{Response ($y$)} & \textbf{$n$, $b$} \\
    \hline
    1 & Sunlight and plant growth & hours of sunlight & growth (cm) & $n=200$,\ $b=0.45$ \\
    2 & Driving and gas use       & miles driven      & gallons used & $n=8$,\ $b=1.2$ \\
    3 & Shoe size and GPA         & shoe size         & GPA & $n=40$,\ $b=0.03$ \\
\end{tabular}
\end{center}

\vspace{0.15in}

\begin{tcolorbox}[colback=white, colframe=black!40,
    title=\textbf{Tier R --- Remediate},
    fonttitle=\bfseries, arc=2mm, left=3mm, right=3mm, top=2mm, bottom=2mm]

\textbf{Scenario 1 only:} $n = 200$,\ $b = 0.45$ (hours of sunlight vs.\ plant growth in cm)

\begin{enumerate}[label=\alph*.]

  \item Interpret the slope $b = 0.45$ in context using the sentence frame:

        \vspace{0.05in}
        ``For each additional \blank{4.5cm}, the predicted \blank{4cm} increases by \blank{2cm}.''

        \vspace{0.1in}

  \item If the true slope $\beta = 0$, would getting $b = 0.45$ from $n = 200$
        observations be surprising? \textbf{Circle one:}

        \vspace{0.05in}
        \hspace{1cm}\textbf{Very surprising} \qquad \textbf{Somewhat surprising} \qquad \textbf{Not surprising}

        \vspace{0.05in}
        Explain briefly: \writelines{2}

  \item What does it mean to say $\beta = 0$ in this context?

        \writelines{2}

\end{enumerate}
\end{tcolorbox}

\begin{tcolorbox}[colback=white, colframe=black!40,
    title=\textbf{Tier A --- Approaching Proficiency},
    fonttitle=\bfseries, arc=2mm, left=3mm, right=3mm, top=2mm, bottom=2mm]

\textbf{All 3 scenarios.}

\begin{enumerate}[label=\alph*.]

  \item Interpret the slope in context for each scenario (one sentence each).

        Scenario 1: \writelines{1}

        Scenario 2: \writelines{1}

        Scenario 3: \writelines{1}

  \item Based on the sample size and slope value, do you think the data provide
        convincing evidence that $\beta \ne 0$? Explain your reasoning (informal --- no
        formal test yet).

        Scenario 1: \writelines{2}

        Scenario 2: \writelines{2}

        Scenario 3: \writelines{2}

  \item For which scenario are you \textbf{most} confident the slope is real?
        Which are you \textbf{least} confident about? Why?

        Most confident: \blank{3cm} \quad Least confident: \blank{3cm}

        Reasoning: \writelines{2}

\end{enumerate}
\end{tcolorbox}

\begin{tcolorbox}[colback=white, colframe=black!40,
    title=\textbf{Tier E --- Extension},
    fonttitle=\bfseries, arc=2mm, left=3mm, right=3mm, top=2mm, bottom=2mm]

Complete all of Tier A above. Then answer:

\begin{enumerate}[label=\alph*., start=4]

  \item Scenarios 1 and 2 have similar slope magnitudes ($b \approx 0.45$ vs.\ $b \approx 1.2$),
        but very different sample sizes ($n = 200$ vs.\ $n = 8$). What does this suggest
        about what matters when deciding whether a slope is real?

        \writelines{3}

  \item Why can't we simply look at whether $b \ne 0$ to conclude that $\beta \ne 0$?

        \writelines{3}

\end{enumerate}
\end{tcolorbox}

\end{document}
